
\preface

Clinical Physiology: A muscle centered approach is a manuscript (book) in progress. The version you are reading was released \today.

\vspace{0.2in}

This project emerged late in my sabbatical year, early April 2022. I had been working on a book for my course sequence titled \textit{Clinical Inquiry} and was unable to break through some theory to practice barriers. I had been meeting for weekly discussions with my friend Bog Sniezek, who is a systems scientist that was at Draper Labs in Cambridge and is now a consultant for the Air Force. During our conversation about the nature and function of unifying systems (Bog's theory), he inspired some thoughts about how physical therapists should learn clinical physiology. It is a foundational knowledge for all areas of practice. The idea that struck me on the basis of unifying systems theory was a muscle-centered approach. From the perspective of a physical therapist, physiology is the generative (muscle cell) and supportive (other systems) role in movement. The tension developed by muscles, which is developed by muscle cells, is at the core of movement. Muscle cells developing tension are a necessary but not sufficient condition for movement. Muscles have always been included in the clinical physiology course that I teach. In fact, there is no new material in the book from the perspective of what is being taught. The only thing new is that all things are considered from the perspective of the muscle fiber. Since my sabbatical was close to completion and I did not have a draft of my original plan (a book on clinical inquiry) and I had a \textit{Clinical Physiology} course to teach in a couple of months; I made the decision to start working on this text. In my 30-year history of teaching this topic, I was never satisfied with the textbook options. It was a Goldilocks situation. I either spent class time explaining to students what they did not need to know and adding topics related to movement to a medical physiology textbook; or I spent time adding clinical topics to an exercise physiology textbook. Either way, it discouraged students from reading the book and focusing on the lectures. However, the lectures were dense and students would often say things such as: \textit{"I wish I had a transcript of the lecture."} This book is a transcript of those lectures. A secondary goal of this book is that students will learn how to learn through graduate level academic reading.

\vspace{0.2in}

A muscle-centered approach places the muscle fiber (cell) at the center of the inquiry. \textbf{Part I} establishes preliminaries to put muscle fiber in context.  \textbf{Part II} focuses on what a muscle fiber does as a system. Muscle fibers create tension, a specific type of force. The process of creating tension is an inward looking process. We consider the cellular, mechanical, and chemical mechanisms of the muscle fiber. \textbf{Part III} focuses on the necessary supporting systems for muscle fiber. Each muscle fiber receives and releases material from its surrounding environment (extra cellular fluid). Several physiological systems are integrated to ensure the stability of that environment. Parts I through III are covered in the author's DPT course, \textit{Clinical Physiology}, at Plymouth State University.

\vspace{0.2in}

\textbf{Part IV} (in progress) will focus on muscle molecular biology. Specifically, how the muscle continues to fulfill its role of creating tension without being a burden on the physiological support systems. Topics include atrophy, hypertrophy, and repair. \textbf{Part V} (in progress) will include a variety of integrated exercises. These topics have been selected to highlight the need for physiological coordination.

\subsection*{Clinical Physiology in a DPT curriculum}
A DPT program has the challenge of educating individuals of diverse backgrounds in the core set of knowledge to become physical therapists. Clinical physiology is a foundational knowledge. The book fits best in a first term course, or a first course on human physiology. The book assumes that the reader has taken prerequisite courses in anatomy \& physiology, physics, and chemistry. It is not an exercise physiology text, though it has a lot of content that would be found in an exercise physiology text. It is not a medical physiology text, though it has a lot of content that would be found in a medical physiology text. The book includes a foundational knowledge of the physiology for a physical therapist. That knowledge is buried in a previously undiscovered textbook space that combines exercise physiology and medical physiology. 

\section*{Tips for the Reader \& Student}
Read actively. Question what you read as you read. Engage in what you read while you read. Every chapter begins with a brief introduction to orient your mind towards what you will be reading about; a mini table of contents so you can see the general layout; and objectives for you to start to think about what you will be learning. It is recommended that you start by reading these and then skim the rest of the chapter as a second step. This process helps you to read actively. During your pre-read skim, you should make mental note of the headings, subheadings, figures, and equations. The purpose of skimming is to get a sense of what is coming up as you start reading the chapter. When you start reading, read as if you were listening to the author speak. If you get a question in your head as you listen to me, write it down. Proceed to read even if you are not completely sure you understand what you are reading. Sometimes, a concept becomes clearer as you proceed. If not, make note of what you are not understanding and bring it to class. Identifying and clearing confusion is one of the purposes of class time; the second purpose of class time is application, depth, and breadth.

Each chapter ends with a summary that includes the \textit{Next Steps} to help you understand what has been covered and where we are headed. At the end of each chapter you will also find sample questions. You should be comfortable answering these questions. But you should also be comfortable manipulating them, turning them into other versions, making them easy or hard multiple choice questions based on the response options. Often the difficulty of a multiple choice question is not the question, the neighbors of the correct response option. Use the review questions to drive discussions during group study time. Yes, study in groups. Group study is the only thing you can do beyond class time to make sure you are not stuck in your own biased trap of understanding. You do not know what you do not know. You need others to bring you from unconscious incompetence to conscious incompetence, the first step towards conscious competence. 

\section*{A Note about Chapter Objectives}

The chapter objectives include specific behavioral terms. There is an ordered relationship between the behavioral terms that represents increasing and nested expectations. 

To \textbf{define} is to know the words being used and to know what they represent. Whether you can define something can be tested with a straightforward question such as identifying the correct definition from a set of definitions; whether you can write the definition; respond true or false to a proposed definition for a term; or identify the correct term from a set of terms when presented with a definition. 

To \textbf{explain} goes beyond defining, but it also implies that you can define the required terms. Explaining is the act associated with understanding. If you understand something, you can explain it. Objectives are typically written as acts - what the person learning can do if they have achieved the objective. Since understanding is observed as an act of explaining, we will use the objective "explain" when there is something that you must "understand". How do you demonstrate that you understand? You explain what it is that you understand. To explain, simultaneously holding a set of definitions together includes their relationship to each other. You can explain a model by demonstrating that you know all the definitions involved and describe how they relate to each other and that you can communicate that explanation. Whether you can explain something can be tested by asking about several aspects of what you are to explain or asking you to infer which part is missing; or by asking about one part and asking you to infer the other part. The point is to explain, which goes beyond defining and has an expectation of definition. But to define does not require you to explain. 

To \textbf{evaluate} requires you to make judgments and infer beyond the thing you are asked to evaluate. To evaluate implies that you can define and explain the thing (or concept) you are asked to evaluate. Being tested for whether you can evaluate can include asking you to infer to or from something that has not been directly covered or discussed, going beyond the thing (or concept) you are asked to evaluate to consider its implications in different situations. 

\textbf{Providing an example} is an objective based on a specific task and is considered a form of asking you to evaluate since it goes beyond a particular thing (or concept) to a totally new thing (or concept). Asking you to provide an example of a model requires you to know the concept of a model well enough that you can use your background knowledge to come up with something you are already familiar with that fits what you know about models. That level of knowledge requires you to evaluate a model because you evaluated your background knowledge and determined that it is a model. If a chapter objective asks you to provide an example of a model that is useful to physical therapy practice, it assumes that you will be able to define, explain, and evaluate the concept of a model and that you have background knowledge about physical therapy practice that you can use to evaluate models used in practice to consider what those models are and whether they are useful. Providing an example is similar to the objective \textbf{Apply}. With both of these objectives, you are asked to use your knowledge of the general concept (or thing) to a specific instance. If the concept was fruit, then an example is an apple. You can also apply the concept of fruit by listing several fruits, including what makes them all fruits.

\textbf{Compare and contrast} is also an objective based on a specific task that requires evaluation of two (or more) things (or concepts) to identify what is similar and different about them. This can occur with the general concept (e.g., fruit) by comparing fruits with vegetables, or with specifics (apples) by comparing different specific fruits (apples and oranges), or by comparing different specifics from each concept (apples and broccoli). Comparison and contrasting are an important higher-level evaluative process that helps you develop a deeper understanding.

\section*{Updates to the Third Version}

The third draft of this manuscript is now released as a **single PDF** as opposed to separate chapters. Therefore, it comes complete with **front matter** (though no back matter). Each chapter now has a **summary** as well as a few (**at least five**) sample questions at the end of the chapter.  

This draft is still not a completed manuscript, but is ready to receive **feedback from colleagues for peer review and possible submission to publishers.  

I had hoped to complete a glossary, an index, and improved contextual links between chapters**, but alas, these are deferred to a future draft.  

\section*{Updates to the Fourth Version}

This fourth version of the manuscript marks another step toward completion. It remains a single PDF, now with both front matter and back matter included.  

There are two new chapters:
\begin{itemize}
  \item The first completes \textbf{Part III: Muscle Support} and focuses on \textbf{digestion, absorption, and metabolism}.
  \item The second is the first chapter in \textbf{Part IV: Integration}, covering \textbf{Fick’s equation}.
\end{itemize}

I have expanded the \textbf{sample questions and added reflection questions} at the end of each chapter to reinforce key concepts. This draft has now undergone \textbf{peer review}, and I have decided to keep it perpetually as an \textbf{Open Educational Resource (OER)}.  

The cost of DPT education is already high, and adding textbook expenses only increases the burden. Moreover, as professors, we are already compensated for our work: There is no need for me to be paid twice for the contributions I am making.  

This is also the first version to include a \textbf{glossary and an index}, both of which will continue to evolve in future versions.  

As with each iteration, I assume full responsibility for any remaining errors. All I can say to the reader is:  
\textit{Mea culpa, veniam peto.}

\vspace{\baselineskip}
\begin{flushright}\noindent
Ashland, New Hampshire\hfill {\it Sean  Collins}\\
May 2025\hfill {\it \hspace{1cm} }\\
\end{flushright}


%\newpage
%\vspace{\baselineskip}
%\begin{flushright}\noindent
%Boscawen, New Hampshire\hfill {\it Sean  Collins}\\
%May 2024\hfill {\it \hspace{1cm} }\\
%\end{flushright}


